\section{Corpus Map Across 82 Study Families}\label{sec:corpus-map}
\subsection{Generating layers, receiving projections and threat polarities}
Table~\ref{tab:core-family-map} maps all \FinalCore{} core families by signal-generating layers, receiving projections and threat polarity. Multi-layer families retain every reported layer in their generating and observing sets.
\begin{table*}[!t]\centering\footnotesize
\caption{Complete map of the 58 core study families. Each family appears once; layer and type labels may be multiple and do not imply one-to-one causal paths. The blocks group source layers, not evidential quality.}\label{tab:core-family-map}
\setlength{\tabcolsep}{2pt}
\begin{minipage}[t]{0.49\textwidth}\vspace{0pt}\centering
\begin{tabularx}{\linewidth}{@{}Ycc@{}}\toprule
Core family & $G\,\to\,O$ & Type\\\midrule
\multicolumn{3}{@{}l}{\textbf{hardware generation, P1 (8)}}\\\midrule
Edge EM leakage~\cite{alhalawani2026physicaledge} & $1\,\to\,1$ & \texttt{I}\\
Energon~\cite{chaudhuri2025energon} & $1\,\to\,1$ & \texttt{II}\\
MoEcho~\cite{ding2025moecho} & $1\,\to\,1$ & \texttt{I}\\
I Know What You Said~\cite{gao2025iknow} & $1\,\to\,1$ & \texttt{I}\\
Behind Bars~\cite{gu2026behindbars} & $1\,\to\,1$ & \texttt{I}\\
PAGE404~\cite{nazaraliyev2026page404} & $1\,\to\,1$ & \texttt{I}\\
Spill the Beans~\cite{spillbeans2025} & $1\,\to\,1$ & \texttt{I}\\
ArchSnoop~\cite{weng2026archsnoop} & $1\,\to\,1$ & \texttt{II}\\
\midrule
\multicolumn{3}{@{}l}{\textbf{hardware and service, P1/P2 (8)}}\\\midrule
SparLeak~\cite{chen2026sparleak} & $1,2\,\to\,1$ & \texttt{I}\\
Shape and Substance~\cite{hadad2026shape} & $1,2\,\to\,1$ & \texttt{I}\\
TDXRay~\cite{hornetz2026tdxray} & $1,2\,\to\,1$ & \texttt{I}\\
SparSEEty~\cite{jo2026sparseety} & $1,2\,\to\,1$ & \texttt{I}\\
OTRO~\cite{lee2026otro} & $1,2\,\to\,1$ & \texttt{I}\\
Detokenization Leaks~\cite{weiss2026detokenization} & $1,2\,\to\,1$ & \texttt{I}\\
LLM-WFIN~\cite{jiao2025llmwfin} & $1,2\,\to\,1,3$ & \texttt{I}\\
LLM-serving RCA~\cite{scheinert2026rca} & $1,2\,\to\,2$ & \texttt{IV}\\
\midrule
\multicolumn{3}{@{}l}{\textbf{service generation, P2 (first 13 of 22)}}\\\midrule
SplitDecoy~\cite{politis2026splitdecoy} & $2\,\to\,2$ & \texttt{I}\\
KVGov~\cite{addagada2026kvgov} & $2\,\to\,2,4$ & \texttt{I}\\
LLMs Have Rhythm~\cite{alhazbi2025rhythm} & $2\,\to\,3$ & \texttt{II}\\
Remote timing~\cite{carlini2024remote} & $2\,\to\,3$ & \texttt{I}\\
NetEcho~\cite{netecho2025} & $2\,\to\,3$ & \texttt{I}\\
Prompt Echo~\cite{prompt_echo2026} & $2\,\to\,3$ & \texttt{I}\\
Speculative decoding~\cite{wei2024speculation} & $2\,\to\,3$ & \texttt{I}\\
Remote keylogging~\cite{weiss2024prompt} & $2\,\to\,3$ & \texttt{I}\\
Whisper Leak~\cite{whisperleak2025} & $2\,\to\,3$ & \texttt{I}\\
Cache contention~\cite{abubakar2026contention} & $2\,\to\,4$ & \texttt{I}\\
Routing metadata~\cite{chen2026routingmetadata} & $2\,\to\,4$ & \texttt{I}\\
SafeKV~\cite{chu2025safekv} & $2\,\to\,4$ & \texttt{I}\\
ML privacy channels~\cite{debenedetti2024privacy} & $2\,\to\,4$ & \texttt{I}\\
\bottomrule\end{tabularx}\end{minipage}\hfill
\begin{minipage}[t]{0.49\textwidth}\vspace{0pt}\centering
\begin{tabularx}{\linewidth}{@{}Ycc@{}}\toprule
Core family & $G\,\to\,O$ & Type\\\midrule
\multicolumn{3}{@{}l}{\textbf{service generation, P2 (remaining 9 of 22)}}\\\midrule
Early Bird~\cite{earlybird2025} & $2\,\to\,4$ & \texttt{I}\\
Prompt-cache audit~\cite{gu2025promptcaching} & $2\,\to\,4$ & \texttt{I}\\
InputSnatch~\cite{inputsnatch2024} & $2\,\to\,4$ & \texttt{I}\\
Leaky Language Models~\cite{majidi2026leakylms} & $2\,\to\,4$ & \texttt{I}\\
PromptPeek~\cite{promptpeek2025} & $2\,\to\,4$ & \texttt{I}\\
SpliceLeak~\cite{spliceleak2026} & $2\,\to\,4$ & \texttt{I}\\
Time Will Tell~\cite{timewilltell2024} & $2\,\to\,4$ & \texttt{I}\\
CachePrune~\cite{wu2026cacheprune} & $2\,\to\,4$ & \texttt{I}\\
MoE prompt stealing~\cite{yona2024stealing} & $2\,\to\,4$ & \texttt{I}\\
\midrule
\multicolumn{3}{@{}l}{\textbf{cross-layer sources (5)}}\\\midrule
ClawGuard~\cite{gan2026clawguard} & $1,5\,\to\,1$ & \texttt{IV}\\
MIRAGE~\cite{revankar2026mirage} & $2,5\,\to\,2$ & \texttt{IV}\\
Cloud traffic benchmark~\cite{pouryousef2026traffic} & $2,5,8\,\to\,3$ & \texttt{I,II}\\
Web-agent unmasking~\cite{fayolle2026webagents} & $3,7\,\to\,3,7$ & \texttt{II,III}\\
LumiMAS~\cite{solomon2026lumimas} & $4,5\,\to\,4,6$ & \texttt{IV}\\
\midrule
\multicolumn{3}{@{}l}{\textbf{orchestration/workflow sources (15)}}\\\midrule
AgentPrint~\cite{agentprint2025} & $5\,\to\,3$ & \texttt{I,II}\\
Cycle detection~\cite{george2025cycles} & $5\,\to\,4,6$ & \texttt{IV}\\
SIGLEAK~\cite{geng2026sigleak} & $5,6\,\to\,4,6$ & \texttt{I}\\
ProcGrep~\cite{oderinwale2026procgrep} & $5,6\,\to\,4,6$ & \texttt{II}\\
AgentProv~\cite{wang2026agentprov} & $5,6\,\to\,6$ & \texttt{II}\\
Lidar~\cite{wang2026lidar} & $5,6\,\to\,6$ & \texttt{II}\\
Dangol workflow~\cite{dangol2026workflow} & $5,6,8\,\to\,3,6,8$ & \texttt{I}\\
Campaign attribution~\cite{park2026crosscampaign} & $5,6,8\,\to\,6,8$ & \texttt{II}\\
WRA~\cite{jeong2026wra} & $5,7\,\to\,3$ & \texttt{I}\\
PersonaFingerprint~\cite{song2026personafingerprint} & $5,7\,\to\,3$ & \texttt{I}\\
Whose Agent?~\cite{kang2026whoseagent} & $5,7\,\to\,7$ & \texttt{II}\\
AgentTell~\cite{shahriar2026agenttell} & $5,7\,\to\,7$ & \texttt{I}\\
UI Traces~\cite{uitraces2026} & $5,7\,\to\,7$ & \texttt{II}\\
FP-Agent~\cite{wang2026fpagent} & $5,7\,\to\,7$ & \texttt{II}\\
AHB~\cite{zhu2026turing} & $5,7\,\to\,7$ & \texttt{III}\\
\bottomrule\end{tabularx}\end{minipage}
\par\smallskip\raggedright $G$ denotes signal-generation layers and $O$ the receiving observation layers. Layer indices: 1 hardware; 2 serving/runtime; 3 transport; 4 application/API; 5 orchestration; 6 tools; 7 browser/environment; 8 multi-agent communication. Types: I confidentiality/privacy; II attribution; III human--agent distinction/evasion; IV authorized diagnostics. These are objectives and observer roles, not mutually exclusive security guarantees. Privileged operator receivers and physical receivers are marked by the layer they inspect. Source versions, locators and normalization of prose-only layer descriptions are in the artifact.
\end{table*}


\paragraph{Cross-layer synthesis}
A remote network observer, a same-site script and a privileged diagnostic hook expose different views even when they share an attribution target. WRA's read-only privacy decoys, AHB's stateful evasion and SplitDecoy's exposed return path illustrate why receiving permissions and protected objectives must accompany the generating layer \cite{jeong2026wra,zhu2026turing,politis2026splitdecoy}.

Layer incidence highlights where the systematization's comparison problem arises. In the 58-family map, signal-generation labels cluster at P2 serving/runtime (32 families) and P5 orchestration (19); receiving views cluster at P4 application/API (18), P1 hardware (16), and P3 transport (14). Generation at P3/P4 is rare (one family each), despite frequent observation at those layers. This source/receiver asymmetry makes the transfer question concrete: protection at a frequent source layer cannot be presumed to cover a different receiving projection. The map records layer sets, so it does not assign a causal edge to each source/receiver pair.

Purpose labels further show why a single “agent side channel” category is insufficient: Type I appears 41 times, Type II 13, Type III 2, and Type IV 5. The 32 map-only families extend coverage beyond mitigation comparisons: 15 list serving origins, 11 hardware and 11 orchestration; 19 concern confidentiality, nine attribution, one evasion and all five authorized-diagnostic families. Mapping these families locates privacy targets, operator capabilities and candidate validation environments even when no scoped mitigation record enters the ledger. Counts are overlapping family-label incidences; the audit then applies its reuse rules to the 39 scoped result records.

\subsection{Relationship to prior systematizations}
Prior surveys map compound-system threats, TEE boundaries and privacy contracts \cite{banerjee2024compound,forough2026secrets,kotevska2026privacyfoundations}. WF already supplies feature leakage analysis, cost trade-offs and environment-shift tests \cite{cai2014systematic,juarez2014critical,juarez2016wtfpad,wang2017walkietalkie}. Mathews et al. re-evaluate nine defenses with stronger attacks and an on-network DynaFlow implementation \cite{mathews2023critical}. Adaptive attacks and online-cost evaluation are established foundations.

\paragraph{What this synthesis adds}
Our unit is a \emph{scoped result}, so execution mode, protected state and cost remain attached when one paper reports several mechanisms. The 39-record census exposes these within-paper substitutions; the 82-family holdout matrix separates evidence for each transfer axis. This connects evidence that earlier taxonomies organize by attack or mechanism: one family can contain several execution modes, a receiving projection can cross system layers, and utility/cost can belong to a different operating point from inference. The synthesis preserves these distinctions when translating a reported result into a deployment decision.
